% GENERATED FILE. Edit canonical lessons, then run npm run generate:books.
% canonical-source-hash: fnv1a64:39ccb41a6eead5bb
% canonical-lessons: TE-C230-kavita, TE-C230-citram, TE-C230-kemera, TE-C230-vinodam, TE-C230-natudu

\chapter{Poem, Painting, Camera, Entertainment, Actor}
\label{ch:te-poem-painting-camera-entertainment-actor}
\hlchaptermodality{230}

\begin{chapteropening}
\textbf{By the end of this chapter:} \emph{I can say a poem, a painting, a camera, entertainment and an actor.}

Complete the last lesson of chapter 230: I can say a poem, a painting, a camera, entertainment and an actor.
\end{chapteropening}

\section[\texorpdfstring{\te{కవిత} (kavita)}{కవిత (kavita)}]{\te{కవిత} (kavita) — a poem}

\label{lesson:TE-C230-kavita}



\begin{quote}
Before the new one: say the Telugu for a raindrop, then the Telugu for fog.
\end{quote}

\subsection*{You'll want to know: \te{కవిత}}
\textbf{\te{కవిత}} — \emph{kavita} — \textquotedblleft{}a poem\textquotedblright{}.

\begin{grammarlens}[title={What you've built}]
One more word for this chapter.
\end{grammarlens}

\subsection*{Your turn}
\begin{itemize}
\raggedright
  \item \emph{Say it:} \emph{kavita}
  \item \emph{Say it:} \emph{kavita}, once more
  \item \emph{Say it:} say \emph{kavita}
  \item \emph{Recall:} say the Telugu for a raindrop, then the Telugu for fog, then say \emph{kavita} again
\end{itemize}

\begin{tcolorbox}[breakable,colback=teal!4,colframe=teal!35!black,title={Before you move on}]
What does \te{కవిత} mean? (\textquotedblleft{}A poem\textquotedblright{}.) Say it once more.
\end{tcolorbox}
\section[\texorpdfstring{\te{చిత్రం} (citraṁ)}{చిత్రం (citraṁ)}]{\te{చిత్రం} (citraṁ) — a painting, a picture}

\label{lesson:TE-C230-citram}



\begin{quote}
Before the new one: say the Telugu for fog, then the Telugu for a poem.
\end{quote}

\subsection*{You'll want to know: \te{చిత్రం}}
\textbf{\te{చిత్రం}} — \emph{citraṁ} — \textquotedblleft{}a painting, a picture\textquotedblright{}.

\begin{grammarlens}[title={What you've built}]
One more word for this chapter.
\end{grammarlens}

\subsection*{Your turn}
\begin{itemize}
\raggedright
  \item \emph{Say it:} \emph{citraṁ}
  \item \emph{Say it:} \emph{citraṁ}, once more
  \item \emph{Say it:} say \emph{citraṁ}
  \item \emph{Recall:} say the Telugu for fog, then the Telugu for a poem, then say \emph{citraṁ} again
\end{itemize}

\begin{tcolorbox}[breakable,colback=teal!4,colframe=teal!35!black,title={Before you move on}]
What does \te{చిత్రం} mean? (\textquotedblleft{}A painting, a picture\textquotedblright{}.) Say it once more.
\end{tcolorbox}
\section[\texorpdfstring{\te{కెమెరా} (kemerā)}{కెమెరా (kemerā)}]{\te{కెమెరా} (kemerā) — a camera}

\label{lesson:TE-C230-kemera}



\begin{quote}
Before the new one: say the Telugu for a poem, then the Telugu for a painting.
\end{quote}

\subsection*{You'll want to know: \te{కెమెరా}}
\textbf{\te{కెమెరా}} — \emph{kemerā} — \textquotedblleft{}a camera\textquotedblright{}.

\begin{grammarlens}[title={What you've built}]
One more word for this chapter.
\end{grammarlens}

\subsection*{Your turn}
\begin{itemize}
\raggedright
  \item \emph{Say it:} \emph{kemerā}
  \item \emph{Say it:} \emph{kemerā}, once more
  \item \emph{Say it:} say \emph{kemerā}
  \item \emph{Recall:} say the Telugu for a poem, then the Telugu for a painting, then say \emph{kemerā} again
\end{itemize}

\begin{tcolorbox}[breakable,colback=teal!4,colframe=teal!35!black,title={Before you move on}]
What does \te{కెమెరా} mean? (\textquotedblleft{}A camera\textquotedblright{}.) Say it once more.
\end{tcolorbox}
\section[\texorpdfstring{\te{వినోదం} (vinōdaṁ)}{వినోదం (vinōdaṁ)}]{\te{వినోదం} (vinōdaṁ) — entertainment}

\label{lesson:TE-C230-vinodam}



\begin{quote}
Before the new one: say the Telugu for a painting, then the Telugu for a camera.
\end{quote}

\subsection*{You'll want to know: \te{వినోదం}}
\textbf{\te{వినోదం}} — \emph{vinōdaṁ} — \textquotedblleft{}entertainment\textquotedblright{}.

\begin{grammarlens}[title={What you've built}]
One more word for this chapter.
\end{grammarlens}

\subsection*{Your turn}
\begin{itemize}
\raggedright
  \item \emph{Say it:} \emph{vinōdaṁ}
  \item \emph{Say it:} \emph{vinōdaṁ}, once more
  \item \emph{Say it:} say \emph{vinōdaṁ}
  \item \emph{Recall:} say the Telugu for a painting, then the Telugu for a camera, then say \emph{vinōdaṁ} again
\end{itemize}

\begin{tcolorbox}[breakable,colback=teal!4,colframe=teal!35!black,title={Before you move on}]
What does \te{వినోదం} mean? (\textquotedblleft{}Entertainment\textquotedblright{}.) Say it once more.
\end{tcolorbox}
\section[\texorpdfstring{\te{నటుడు} (naṭuḍu)}{నటుడు (naṭuḍu)}]{\te{నటుడు} (naṭuḍu) — an actor}

\label{lesson:TE-C230-natudu}



\begin{quote}
Before the new one: say the Telugu for a camera, then the Telugu for entertainment.
\end{quote}

\subsection*{You'll want to know: \te{నటుడు}}
\textbf{\te{నటుడు}} — \emph{naṭuḍu} — \textquotedblleft{}an actor\textquotedblright{}.

\begin{grammarlens}[title={What you've built}]
One more word for this chapter.
\end{grammarlens}

\subsection*{Your turn}
\begin{itemize}
\raggedright
  \item \emph{Say it:} \emph{naṭuḍu}
  \item \emph{Say it:} \emph{naṭuḍu}, once more
  \item \emph{Say it:} say \emph{naṭuḍu}
  \item \emph{Recall:} say the Telugu for a camera, then the Telugu for entertainment, then say \emph{naṭuḍu} again
\end{itemize}

\begin{tcolorbox}[breakable,colback=teal!4,colframe=teal!35!black,title={Before you move on}]
What does \te{నటుడు} mean? (\textquotedblleft{}An actor\textquotedblright{}.) Say it once more.
\end{tcolorbox}
